% Options for packages loaded elsewhere
% Options for packages loaded elsewhere
\PassOptionsToPackage{unicode}{hyperref}
\PassOptionsToPackage{hyphens}{url}
\PassOptionsToPackage{dvipsnames,svgnames,x11names}{xcolor}
%
\documentclass[
  english,
  11pt,
]{article}
\usepackage{xcolor}
\usepackage[margin=25mm]{geometry}
\usepackage{amsmath,amssymb}
\setcounter{secnumdepth}{-\maxdimen} % remove section numbering
\usepackage{iftex}
\ifPDFTeX
  \usepackage[T1]{fontenc}
  \usepackage[utf8]{inputenc}
  \usepackage{textcomp} % provide euro and other symbols
\else % if luatex or xetex
  \usepackage{unicode-math} % this also loads fontspec
  \defaultfontfeatures{Scale=MatchLowercase}
  \defaultfontfeatures[\rmfamily]{Ligatures=TeX,Scale=1}
\fi
\usepackage{lmodern}
\ifPDFTeX\else
  % xetex/luatex font selection
\fi
% Use upquote if available, for straight quotes in verbatim environments
\IfFileExists{upquote.sty}{\usepackage{upquote}}{}
\IfFileExists{microtype.sty}{% use microtype if available
  \usepackage[]{microtype}
  \UseMicrotypeSet[protrusion]{basicmath} % disable protrusion for tt fonts
}{}
\makeatletter
\@ifundefined{KOMAClassName}{% if non-KOMA class
  \IfFileExists{parskip.sty}{%
    \usepackage{parskip}
  }{% else
    \setlength{\parindent}{0pt}
    \setlength{\parskip}{6pt plus 2pt minus 1pt}}
}{% if KOMA class
  \KOMAoptions{parskip=half}}
\makeatother
% Make \paragraph and \subparagraph free-standing
\makeatletter
\ifx\paragraph\undefined\else
  \let\oldparagraph\paragraph
  \renewcommand{\paragraph}{
    \@ifstar
      \xxxParagraphStar
      \xxxParagraphNoStar
  }
  \newcommand{\xxxParagraphStar}[1]{\oldparagraph*{#1}\mbox{}}
  \newcommand{\xxxParagraphNoStar}[1]{\oldparagraph{#1}\mbox{}}
\fi
\ifx\subparagraph\undefined\else
  \let\oldsubparagraph\subparagraph
  \renewcommand{\subparagraph}{
    \@ifstar
      \xxxSubParagraphStar
      \xxxSubParagraphNoStar
  }
  \newcommand{\xxxSubParagraphStar}[1]{\oldsubparagraph*{#1}\mbox{}}
  \newcommand{\xxxSubParagraphNoStar}[1]{\oldsubparagraph{#1}\mbox{}}
\fi
\makeatother


\usepackage{longtable,booktabs,array}
\newcounter{none} % for unnumbered tables
\usepackage{calc} % for calculating minipage widths
% Correct order of tables after \paragraph or \subparagraph
\usepackage{etoolbox}
\makeatletter
\patchcmd\longtable{\par}{\if@noskipsec\mbox{}\fi\par}{}{}
\makeatother
% Allow footnotes in longtable head/foot
\IfFileExists{footnotehyper.sty}{\usepackage{footnotehyper}}{\usepackage{footnote}}
\makesavenoteenv{longtable}
\usepackage{graphicx}
\makeatletter
\newsavebox\pandoc@box
\newcommand*\pandocbounded[1]{% scales image to fit in text height/width
  \sbox\pandoc@box{#1}%
  \Gscale@div\@tempa{\textheight}{\dimexpr\ht\pandoc@box+\dp\pandoc@box\relax}%
  \Gscale@div\@tempb{\linewidth}{\wd\pandoc@box}%
  \ifdim\@tempb\p@<\@tempa\p@\let\@tempa\@tempb\fi% select the smaller of both
  \ifdim\@tempa\p@<\p@\scalebox{\@tempa}{\usebox\pandoc@box}%
  \else\usebox{\pandoc@box}%
  \fi%
}
% Set default figure placement to htbp
\def\fps@figure{htbp}
\makeatother



\ifLuaTeX
\usepackage[bidi=basic,shorthands=off]{babel}
\else
\usepackage[bidi=default,shorthands=off]{babel}
\fi
\ifLuaTeX
  \usepackage{selnolig} % disable illegal ligatures
\fi


\setlength{\emergencystretch}{3em} % prevent overfull lines

\providecommand{\tightlist}{%
  \setlength{\itemsep}{0pt}\setlength{\parskip}{0pt}}



 


\makeatletter
\@ifpackageloaded{caption}{}{\usepackage{caption}}
\AtBeginDocument{%
\ifdefined\contentsname
  \renewcommand*\contentsname{Table of contents}
\else
  \newcommand\contentsname{Table of contents}
\fi
\ifdefined\listfigurename
  \renewcommand*\listfigurename{List of Figures}
\else
  \newcommand\listfigurename{List of Figures}
\fi
\ifdefined\listtablename
  \renewcommand*\listtablename{List of Tables}
\else
  \newcommand\listtablename{List of Tables}
\fi
\ifdefined\figurename
  \renewcommand*\figurename{Figure}
\else
  \newcommand\figurename{Figure}
\fi
\ifdefined\tablename
  \renewcommand*\tablename{Table}
\else
  \newcommand\tablename{Table}
\fi
}
\@ifpackageloaded{float}{}{\usepackage{float}}
\floatstyle{ruled}
\@ifundefined{c@chapter}{\newfloat{codelisting}{h}{lop}}{\newfloat{codelisting}{h}{lop}[chapter]}
\floatname{codelisting}{Listing}
\newcommand*\listoflistings{\listof{codelisting}{List of Listings}}
\makeatother
\makeatletter
\makeatother
\makeatletter
\@ifpackageloaded{caption}{}{\usepackage{caption}}
\@ifpackageloaded{subcaption}{}{\usepackage{subcaption}}
\makeatother
\usepackage{bookmark}
\IfFileExists{xurl.sty}{\usepackage{xurl}}{} % add URL line breaks if available
\urlstyle{same}
% fallback for those not using the hyperref driver hyperxmp:
\makeatletter
\@ifundefined{xmpquote}{\newcommand{\xmpquote}[1]{#1}}{}
\makeatother
\hypersetup{
  pdftitle={How can speed rise while energy falls?},
  pdflang={en},
  colorlinks=true,
  linkcolor={blue},
  filecolor={Maroon},
  citecolor={Blue},
  urlcolor={Blue},
  pdfcreator={LaTeX via pandoc}}


\title{How can speed rise while energy falls?}
\usepackage{etoolbox}
\makeatletter
\providecommand{\subtitle}[1]{% add subtitle to \maketitle
  \apptocmd{\@title}{\par {\large #1 \par}}{}{}
}
\makeatother
\subtitle{A shrinking region changes the energy calculation}
\author{}
\date{2026-09-12}
\begin{document}
\maketitle


\href{../paper-map/article.html}{Start with the paper map} · Main route,
article 1 of 9

If a fluid becomes faster and faster, surely its energy must grow
without limit. That sounds reasonable until we ask a second question:
\textbf{how much fluid is moving at that speed?}

A speeding truck carries much more kinetic energy than a speeding grain
of sand. Speed matters, but so does mass. The proposed vortex
concentrates its fastest motion into a shrinking region. This makes it
possible for the largest speed to rise while the energy in that region
falls.

We will use one familiar energy formula and the volume of a box. Those
are enough to understand this first building block.

\subsection{Start with a small amount of moving
material}\label{start-with-a-small-amount-of-moving-material}

The kinetic energy of a mass \(m\) moving at speed \(v\) is

\[
E=\frac12 mv^2.
\]

Doubling the speed multiplies the energy by four, \textbf{if the mass
stays the same}. But suppose the amount of material moving at that speed
becomes one eighth as large. Then the energy changes by a factor of
\(4/8\), which is one half.

{\def\LTcaptype{none} % do not increment counter
\begin{longtable}[]{@{}lll@{}}
\toprule\noalign{}
Quantity & First stage & Next stage \\
\midrule\noalign{}
\endhead
\bottomrule\noalign{}
\endlastfoot
Mass in the fast region & 1 & 1/8 \\
Typical speed there & 1 & 2 \\
Energy in that region & 1/2 & 1/4 \\
\end{longtable}
}

These are arbitrary consistent units for an arithmetic example. We are
comparing the contents of a region at two stages, not claiming that one
material parcel loses seven eighths of its mass.

Repeat the pattern. The speed becomes 4, then 8, then 16. The energy
becomes \(1/8\), then \(1/16\), then \(1/32\). There is no
contradiction: increasingly fast motion is carried by an increasingly
small amount of material.

\begin{figure}

\centering{

\pandocbounded{\includegraphics[keepaspectratio,alt={A sequence compares rising relative speeds 1, 2, 4 and 8 with falling relative energies 1, one half, one quarter and one eighth.}]{figures/reader.png}}

}

\caption{\label{fig-energy}In this simple example, halving three
dimensions reduces the volume eightfold. Doubling the speed then leaves
half as much energy in the region. Repeating the operation gives the
bars on the right.}

\end{figure}%

\subsection{Where the shrinking mass comes
from}\label{where-the-shrinking-mass-comes-from}

For a fluid of fixed density, the mass in a region is proportional to
its volume. A box half as wide, half as deep, and half as high has one
eighth of the original volume.

A cylinder follows the same counting rule when its radius and height are
both halved: its circular area becomes one quarter as large, and its
height becomes one half as large. Multiply them to get one eighth.

The paper's intense core is approximately a narrowing column. Its radius
and axial length both shrink. They do so at slightly different rates,
but the basic competition is the one we have just calculated:

\textbf{energy in the core depends on shrinking volume times growing
speed squared.}

With the paper's chosen rates, the shrinking volume wins. The core's
energy tends to zero, even as the characteristic speed grows. The entire
flow has other regions, so controlling total energy requires an
additional argument. The core calculation explains compatibility; it
does not replace that argument.
\href{https://cdn.openai.com/pdf/32d9f210-8b73-45e0-91bc-82a30aef8a9a/navier-stokes.pdf}{Paper,
Section 2.1 and Section 10.3}

\subsection{A high peak is different from a large
total}\label{a-high-peak-is-different-from-a-large-total}

Imagine a crowd viewed from above. ``How crowded is the busiest square
metre?'' and ``How many people are in the whole square?'' are different
questions.

Likewise, the largest fluid speed is a local measure. Total energy adds
contributions throughout space. A very high velocity peak confined to a
sufficiently small region need not dominate that sum.

This is why an energy bound alone does not automatically provide a bound
on the largest speed. A total can stay manageable while one increasingly
narrow feature becomes extreme.

The accompanying \href{animations/core_collapse.html}{animation} shows a
smooth illustrative bump getting taller and narrower. In a view that
continually zooms in and adjusts the velocity scale, its normalized
shape stays the same. That is the idea of \textbf{self-similarity}: a
changing physical pattern can have a stable shape when measured with
changing rulers.

The animation is a model of concentration. Its arrows and colours do not
track the trajectories of individual molecules.

\subsection{Does a shrinking region violate
incompressibility?}\label{does-a-shrinking-region-violate-incompressibility}

An incompressible material parcel keeps its volume as it moves. That
seems to clash with the shrinking core, until we distinguish the parcel
from the region.

Think of the bright spot made by a moving spotlight. The spot is a
region defined by a property---being brightly lit. It does not always
contain the same people.

The vortex core is similarly a region defined by intense motion and the
construction's changing scales. Fluid can enter and leave it. Its
shrinking volume therefore does not mean that the same parcel is being
compressed into less space.

In the central flow described by the paper, radial inflow is accompanied
by axial outflow. The full geometry is more detailed, but the accounting
principle remains: fluid does not accumulate merely because motion
points inward in one direction.

\subsection{Where this fits in the bigger
picture}\label{where-this-fits-in-the-bigger-picture}

\textbf{This block removes an apparent energy contradiction. It does not
yet produce a solution of the fluid equations.} A chosen shape can have
the right speed and energy behaviour while requiring an unacceptable
force.

The paper must next build a smooth vortex profile and make its momentum
balance work. Continue with \href{../d6_inner_profile/article.html}{what
happens at the centre}. For a separate calculation of actual length
scales under an assumed reference size, see
\href{../d1_core_size/article.html}{core size versus speed}.

The \href{technical-article.html}{technical article},
\href{../../code/d2_finite_energy/study.ipynb}{notebook}, and
\href{source-notes.md}{source notes} contain the exact scaling formulas,
computational checks, and animation conventions.




\end{document}
